\section{Introduction}
This study evaluates alternatives to document scanning within selected clusters. The baseline
uses the compact representation and scoring approach described in Exa's public vector-database
article~\cite{exa}. A complete scan scores every selected document and returns the best results.
The alternatives use information in the query to select candidates before full scoring. Their
potential benefit depends on whether candidate selection costs less than the scoring it avoids,
while retaining the required results.

Three methods are compared on the same binary document embeddings. Method A scans the documents
in selected clusters. Method B searches those clusters through bitplane splits. Method C looks
up bit patterns in increasing mismatch cost, beginning with the pattern preferred by the query.
C is referred to as \emph{weighted key lookup}; it implements the earlier proposal to start
with the ideal sign pattern and consider progressively less favorable patterns.

The objective is to minimize query latency for a fixed document count, RAM budget, and recall
requirement. Recall is the fraction of correct reference results recovered by a method. The
comparison accounts for both computational work and memory access, since avoiding a document
score can introduce work elsewhere in the query.

The results depend on the query and document distribution. Neither alternative established a
reliable speed advantage over a tuned scan on the evaluated real embeddings. Controlled queries
with fixed important coordinates produced a small key-lookup advantage over a closely related
directly routed scan. When the important coordinates varied by query, bitplane search showed
a large advantage in the constructed case. These findings identify conditions under which
each method is useful; they do not establish a universal winner.

The analysis begins with one query and six documents, followed by the work and storage costs
of each method. Subsequent sections derive recall under explicit assumptions, examine parameter
selection, and report measurements. The final section considers a billion-document collection.
Those latency estimates are forecasts; complete C++ indexes were tested through eight million
documents.

\subsection{A query example}
We use the binary document representation and floating-point query described in Exa's
post~\cite{exa}. In our examples, a stored bit of 1 means $+1$, and a bit of 0 means $-1$. We will use five coordinates so we can follow
the calculation on the page. Let the query be
\[
q=[0.7,\ 0.5,\ -0.4,\ -0.6,\ 0.3].
\]
With that same dot-product score, signs \texttt{11001} make every contribution positive, giving
$0.7+0.5+0.4+0.6+0.3=2.5$. This is the pattern the query would prefer. The following example compares this pattern with the stored documents.

\begin{center}
\begin{tabular}{@{}rcrr@{}}\toprule
Document ID & Stored signs & First-bit cluster & Full score\\\midrule
0 & \texttt{01001} & 0 & 1.1\\
1 & \texttt{00100} & 0 & $-1.3$\\
2 & \texttt{10100} & 1 & 0.1\\
3 & \texttt{11001} & 1 & 2.5\\
4 & \texttt{11000} & 1 & 1.9\\
5 & \texttt{00011} & 0 & $-1.1$\\\bottomrule
\end{tabular}
\end{center}
Here the correct top two are documents 3 and 4. If scores tie, we put the smaller document ID
first, and then use that same rule when checking every method. We can also run this example
against the C++ implementations with \path{examples/small_search.py}.

At coordinate 1, a matching sign contributes $+0.7$ and a mismatch contributes $-0.7$,
reducing the score by $1.4$. At coordinate 5, the reduction is only $0.6$. Query magnitude
therefore determines the cost of an incorrect sign and provides a basis for ordering the search.
For $d$ coordinates, the ideal score is $Q=\sum_{i=1}^{d}|q_i|$. Let $b_i$ denote the decoded
sign ($+1$ or $-1$). If mismatching coordinates have total magnitude $p$, then
\begin{equation}
S(q,b)=Q-2p,
\qquad p=\sum_{i:\,b_i\ne\operatorname{sign}(q_i)}|q_i|.
\label{eq:score}
\end{equation}
Document 4 mismatches only coordinate 5 and therefore has score $S=2.5-2(0.3)=1.9$.
Ranking by decreasing score is equivalent to ranking by increasing mismatch penalty.
Ordinary Hamming distance counts all mismatches equally; the score used here weights them
by the query magnitudes.

\subsection{Shared pipeline}
Suppose the collection contains $N$ documents. Before any queries arrive, we embed the text,
choose a representation width $d$, pack the document signs, and build the index. We also assign
each document to one cluster. A cluster is a stored subset of the collection, and we can form
those subsets by learning from the vectors or by using a fixed sign prefix.

For shortening embeddings, we also refer to Exa's use of Matryoshka training~\cite{exa,mrl}.
The representation is held fixed in this study. One distinction matters for the proposed search:
a useful prefix is not the same thing as the largest query coordinates. Matryoshka does not
sort each query by magnitude, and it does not guarantee that its first 24 sign bits preserve
its nearest neighbors. Shortening the vector and taking its signs are two separate changes;
we hold the resulting representation fixed across the methods we compare.

Once a query arrives, we embed and normalize it, select clusters, and search their documents.
Opening one cluster is one \emph{probe}. With $C$ clusters and $P$ probes, balanced clusters
would give us about $NP/C$ selected documents. If the clusters are not balanced, we need to
add their actual sizes instead. If the objective is to search the whole collection, we can set
$C=P=1$.

\fig{shared}{The common query flow. The three local search methods share the same query and selected clusters when we isolate their work. Longer-vector reranking is optional; it is excluded from the primary binary-score comparison.}{.68}

A production pipeline can also filter documents or rerank candidates with longer vectors or
another model. If those stages are enabled, they belong in the complete request measurement.
For the primary index comparison here, we use cached query embeddings, the same binary score,
and no float reranker. The later text-query evaluation measures complete request time.

The baseline therefore follows Exa's public description of compact scoring inside selected
clusters~\cite{exa}. The measurements here are of our local implementation, not Exa's production
system. We also try two ways of forming clusters: learning groups from similar vectors,
known as IVF, and assigning documents by their first few sign bits. With a $k$-bit prefix,
there are up to $2^k$ possible patterns, not $k$ clusters. B and C change how we search the
selected documents; they do not introduce the shared embedding or quantization techniques.